% figures/dynamics/track_vehicle_relation.tex

\begin{figure}[htbp]
\centering

\begin{tikzpicture}[
node distance=1.5cm,
box/.style={
draw,
rounded corners,
minimum width=4cm,
minimum height=0.9cm,
align=center
}
]

\node[box] (pose3)
{Railway State\\$\mathrm{pose3}$};

\node[box, below=of pose3] (track)
{Track State};

\node[box, below=of track] (vehicle)
{Vehicle State};

\node[box, below left=1.3cm and 2.2cm of vehicle] (wheel)
{Wheelset};

\node[box, below=1.3cm of vehicle] (bogie)
{Bogie};

\node[box, below right=1.3cm and 2.2cm of vehicle] (cg)
{CG State};

\draw[->, thick] (pose3) -- (track);
\draw[->, thick] (track) -- (vehicle);

\draw[->, thick] (vehicle) -- (wheel);
\draw[->, thick] (vehicle) -- (bogie);
\draw[->, thick] (vehicle) -- (cg);

\end{tikzpicture}

\caption{
Vehicle-space hierarchy.
Vehicle states are derived from railway states and form the basis
for wheelset, bogie, and center-of-gravity interpretation.
}

\label{fig:track-vehicle-relation}

\end{figure}
